\subsection*{Choosing a balanced party}
For a first game, combine two front-line fighters with a cleric and a mage behind them. The cleric brings healing and food; the mage brings attacks and protections. A thief or a multi-class character adds lockpicks without necessarily giving up another role. Carry ranged weapons for the rear rank.

\begin{reference}{Class}{Role in the party}
Fighter & Weapons and armour specialist; a good choice for taking attacks from the front.\\
Ranger & A combatant sharing many equipment options with the fighter. Do not assume a ranger has a mage's spell book.\\
Paladin & A good-aligned combatant with Lay on Hands and access to a selection of prayers. Alignment restricts the party's composition.\\
Cleric & Healing, food, protections and help against undead. Can wear armour, but has a more restricted choice of weapons.\\
Mage & Learns spells from a book and scrolls. Physical weapons and protection are limited: keep the mage in the rear rank.\\
Thief & Uses lockpicks on locks. Not every lock can be opened this way; keep keys as well.\\
Multiple classes & Combines the roles of the listed classes. Experience is shared between them, so progress differs from a single-class character.\\
\end{reference}

\textbf{Races.} Humans, elves, half-elves, dwarves, gnomes and halflings offer different class combinations and ability limits. Character generation shows the actual choices available for each selection. An elf cannot receive this version's Raise Dead spell; consider this before making the party. Sex and portrait define the character without changing the controls.

\begin{reference}{Ability}{What to look for}
Strength (STR) & Important for physical combat and forcing doors. Some fighters also display exceptional strength alongside 18.\\
Intelligence (INT) & The mage's ability; consider it when choosing a character who will study magic.\\
Wisdom (WIS) & The cleric's ability; relevant to preparing prayers.\\
Dexterity (DEX) & Helps protect the character. Check its effect on AC, especially with little armour.\\
Constitution (CON) & Important for endurance and hit points. Do not neglect the front rank's HP.\\
Charisma (CHA) & Part of the character sheet and some class requirements. It is not a negotiation command and does not replace dialogue choices.\\
\end{reference}

\textbf{Alignment.} Describes the character's moral outlook and restricts some classes and party combinations. It does not add an ability button. Review the characters before pressing PLAY; you need not accept the first roll.
